\chapter{Experimental Setup}
\section{CERN}
International laboratory in Europe, founded 1959?, maybe little significant history.  

\section{LHC}
All the good stats on the LHC: circumference, (place,) dipole design, magnets, magnetic field.  Design energy and luminosity.  
And pictures!  diagram of LHC, diagram of magnets?  just going through prelim diagrams here!

\section{CMS}
Important stats on CMS: height, weight, etc.  diagram of detector.  

Current status on detector, data-taking, event display -- here? or after everything explained?  

Figure: Event display, Fig. \ref{fig:EventDisplay}

 \begin{figure}[htb]
  \begin{center}
    \includegraphics[width=360pt]{CMS-BW.pdf}
  \end{center}
  \caption[Three-dimensional display of a \Zee event in the CMS detector]{Three-dimensional display of a \Zee event in the CMS detector.}
  \label{fig:EventDisplay}
 \end{figure}

\subsection{Detection of Particle Interactions}

\subsection{Coordinate System}


\subsection{Magnet}
Stats about the magnet!

\subsection{Luminosity}
Figure: Luminosity plot, Fig. \ref{fig:LuminosityVsTime}

 \begin{figure}[htb]
  \begin{center}
    \includegraphics[width=360pt]{CMS-BW.pdf}
  \end{center}
  \caption[Luminosity as a function of time]{Luminosity as a function of time.}
  \label{fig:LuminosityVsTime}
 \end{figure}

\subsection{Calorimeters}
\subsubsection{ECAL}
\subsubsection{HCAL}
\subsection{Muon Systems}
\subsection{Tracker}
\subsection{Pixels}
\subsection{Trigger}
CMS trigger identifies potentially interesting events to save.  
Hardware-based Level-1 Trigger reduces event rate by 
factor of X to Y Hz, computer farm High-Level Trigger reduces rate further to Z Hz.  

Picture: Overall diagram of trigger
\subsubsection{Level-1 Trigger}
General L1 information, specific electron algorithms in online selection section.  

L1 uses FPGA-based algorithms to roughly reconstruct and identify possibly interesting physics objects.  
Parallel chains handle muon and calorimetric (?) objects.  
Subsections on that?  Need to talk about the muon parts, as well as the multi-level calo part.  
The L1 Global Trigger combines information from both chains and makes a single decision based on the trigger menu implemented.  
(Ooh boy, go back to RT2010 stuff!)  

\subsubsection{High-Level Trigger}
General HLT information, specific electron algorithms in online selection section.  

HLT uses version of the standard offline physics-object reconstruction algorithms (BEFORE THIS SECTION?????) 
optimized for fast performance.  
Rate at this stage lower, so more time can be afforded to spend on decision, 
more detailed reconstruction can be done.  


